\SourcePage{109}{112}
\section{Angular momentum}

In deriving momentum conservation in \secref{15}, we used the homogeneity of space
for a closed system of particles. Besides being homogeneous, space is also
isotropic: every direction in it is equivalent. The Hamiltonian of a closed
system must therefore remain unchanged when the system as a whole is rotated
through an arbitrary angle about an arbitrary axis. It is enough to impose
this requirement for an arbitrary infinitesimal rotation.

Let $\delta\boldsymbol\varphi$ be the vector of an infinitesimal rotation. Its
magnitude is the angle of rotation $\delta\varphi$, and its direction is the
axis about which the rotation is made. Under such a rotation, the changes
$\delta\mathbf r_a$ of the particles' radius vectors $\mathbf r_a$ are
\[
  \delta\mathbf r_a=\delta\boldsymbol\varphi\times\mathbf r_a.
\]
An arbitrary function $\psi(\mathbf r_1,\mathbf r_2,\ldots)$ is transformed
into
\[
  \begin{aligned}
  \psi(\mathbf r_1+\delta\mathbf r_1,
       \mathbf r_2+\delta\mathbf r_2,\ldots)
    &=\psi(\mathbf r_1,\mathbf r_2,\ldots)
      +\sum_a\delta\mathbf r_a\boldsymbol\nabla_a\psi={}\\
    &=\psi(\mathbf r_1,\mathbf r_2,\ldots)
      +\sum_a(\delta\boldsymbol\varphi\times\mathbf r_a)
      \!\cdot\boldsymbol\nabla_a\psi={}\\
    &=\left(1+\delta\boldsymbol\varphi\!\cdot
      \sum_a(\mathbf r_a\times\boldsymbol\nabla_a)\right)
      \psi(\mathbf r_1,\mathbf r_2,\ldots).
  \end{aligned}
\]
The expression
\[
  1+\delta\boldsymbol\varphi\!\cdot
  \sum_a(\mathbf r_a\times\boldsymbol\nabla_a)
\]
is the operator of an infinitesimal rotation. The fact that such a rotation
does not change the system's Hamiltonian is expressed (cf. \secref{15}) by the
commutation of the rotation operator with $\hat H$. Since
$\delta\boldsymbol\varphi$ is a constant vector, this%
\SourcePage{110}{113}
\relax\ condition reduces to
\begin{equation}
  \left(\sum_a(\mathbf r_a\times\boldsymbol\nabla_a)\right)\hat H
  -\hat H\left(\sum_a(\mathbf r_a\times\boldsymbol\nabla_a)\right)=0,
  \tag{26.1}
\end{equation}
which expresses a conservation law.

The quantity conserved in a closed system as a consequence of the isotropy
of space is the system's \emph{angular momentum} (cf. Vol.~I, \S~9). Thus,
apart from a constant factor, the operator
$\sum_a(\mathbf r_a\times\boldsymbol\nabla_a)$ must correspond to the total angular
momentum of the system's motion, while each term
$(\mathbf r_a\times\boldsymbol\nabla_a)$ must correspond to the angular momentum of
one particle.

The proportionality factor must be chosen as $-i\hbar$. The resulting
particle angular-momentum operator
$-i\hbar(\mathbf r\times\boldsymbol\nabla)=
\mathbf r\times\hat{\mathbf p}$ then agrees
exactly with the classical expression $\mathbf r\times\mathbf p$. Henceforth we
shall always measure angular momentum in units of $\hbar$. We denote the
operator of the angular momentum of one particle, defined in this way, by
$\hat{\mathbf l}$, and that of the entire system by $\hat{\mathbf L}$. Thus
the particle angular-momentum operator is
\begin{equation}
  \hbar\hat{\mathbf l}=\mathbf r\times\hat{\mathbf p}
  =-i\hbar(\mathbf r\times\boldsymbol\nabla)
  \tag{26.2}
\end{equation}
or, in components,
\[
  \hbar\hat l_x=y\hat p_z-z\hat p_y,\qquad
  \hbar\hat l_y=z\hat p_x-x\hat p_z,\qquad
  \hbar\hat l_z=x\hat p_y-y\hat p_x.
\]

For a system in an external field, angular momentum is not generally
conserved. It may nevertheless be conserved when the field has a suitable
symmetry. In a centrally symmetric field, for example, all spatial directions
from the center are equivalent, so angular momentum about that center is
conserved. When the field has axial symmetry, the same reasoning shows that
what is conserved is the angular-momentum component directed along the symmetry
axis. All these conservation laws of
classical mechanics remain valid in quantum mechanics.

If a system's angular momentum is not conserved, it has no definite value in
a stationary state. In such cases the mean angular momentum in a given
stationary state can sometimes be of interest. Its mean value is readily seen
to vanish in every nondegenerate stationary state. Indeed, the energy is
unchanged when the sign of time is reversed; and since only one stationary
state corresponds to the given energy level, replacing $t$ by $-t$ must leave
the state of the system unchanged%
\SourcePage{111}{114}
\relax. The mean values of all quantities, including angular momentum, must
therefore also remain unchanged. But angular momentum reverses sign under
time reversal, which would give
$\bar{\mathbf L}=-\bar{\mathbf L}$; hence $\bar{\mathbf L}=0$. A direct
calculation confirms this: the integral
$\int\psi^*\hat{\mathbf L}\psi\,dq$, which by definition gives the mean
$\bar{\mathbf L}$, can be evaluated using the fact that nondegenerate wave
functions are real (see the end of \secref{18}). One then finds that
\[
  \bar{\mathbf L}=-i\hbar\int\psi^*
  \left(\sum_a(\mathbf r_a\times\boldsymbol\nabla_a)\right)\psi\,dq
\]
is purely imaginary. Since $\bar{\mathbf L}$ must of course be real, it follows
that $\bar{\mathbf L}=0$.

We now determine the commutation rules between the angular-momentum operators
and the coordinate and momentum operators. From (16.2) we readily obtain
\begin{equation}
  \begin{aligned}
  \{\hat l_x,x\}&=0,&
  \{\hat l_x,y\}&=iz,&
  \{\hat l_x,z\}&=-iy,\\
  \{\hat l_y,y\}&=0,&
  \{\hat l_y,z\}&=ix,&
  \{\hat l_y,x\}&=-iz,\\
  \{\hat l_z,z\}&=0,&
  \{\hat l_z,x\}&=iy,&
  \{\hat l_z,y\}&=-ix.
  \end{aligned}
  \tag{26.3}
\end{equation}
For example,
\[
  \hat l_xy-y\hat l_x
  =\frac1\hbar(y\hat p_z-z\hat p_y)y
   -y(y\hat p_z-z\hat p_y)\frac1\hbar
  =-\frac z\hbar\{\hat p_y,y\}=iz.
\]

All the relations (26.3) may be written in tensor form as
\begin{equation}
  \{\hat l_i,x_k\}=ie_{ikl}x_l,
  \tag{26.4}
\end{equation}
where $e_{ikl}$ is the antisymmetric unit tensor of third
rank\footnote{The antisymmetric unit tensor of third rank $e_{ikl}$ (also
called the unit axial tensor) is defined to be antisymmetric in all three
indices, with $e_{123}=1$. Evidently, of its 27 components, only the 6 for
which $i,k,l$ form a permutation of the numbers 1, 2, 3 are nonzero. A
component is $+1$ if the permutation $i,k,l$ is obtained from 1, 2, 3 by an
even number of pairwise interchanges (transpositions), and is $-1$ for an odd
number. Clearly,
\[
  e_{ikl}e_{ikm}=2\delta_{lm},\qquad e_{ikl}e_{ikl}=6.
\]
The components of the vector $\mathbf C=\mathbf A\times\mathbf B$, the vector
product of $\mathbf A$ and $\mathbf B$, can be expressed through $e_{ikl}$ as
\[
  C_i=e_{ikl}A_kB_l.
\]}
, with summation understood over each pair of repeated ``dummy'' indices.
\SourcePage{112}{115}
The angular-momentum and momentum operators satisfy analogous commutation
relations:
\begin{equation}
  \{\hat l_i,\hat p_k\}=ie_{ikl}\hat p_l.
  \tag{26.5}
\end{equation}
These formulas make it easy to obtain the commutation rules among the
operators for the components of angular momentum. We have
\[
  \begin{aligned}
  \hbar(\hat l_x\hat l_y-\hat l_y\hat l_x)
    &=\hat l_x(z\hat p_x-x\hat p_z)
      -(z\hat p_x-x\hat p_z)\hat l_x={}\\
    &=(\hat l_xz-z\hat l_x)\hat p_x
      -x(\hat l_x\hat p_z-\hat p_z\hat l_x)
      =-iy\hat p_x+ix\hat p_y=i\hbar\hat l_z.
  \end{aligned}
\]
Thus
\begin{equation}
  \{\hat l_y,\hat l_z\}=i\hat l_x,\qquad
  \{\hat l_z,\hat l_x\}=i\hat l_y,\qquad
  \{\hat l_x,\hat l_y\}=i\hat l_z
  \tag{26.6}
\end{equation}
or
\begin{equation}
  \{\hat l_i,\hat l_k\}=ie_{ikl}\hat l_l.
  \tag{26.7}
\end{equation}

Exactly the same relations hold for the operators
$\hat L_x$, $\hat L_y$, and $\hat L_z$ of the system's total angular
momentum. Indeed, since the angular-momentum operators of different particles
commute with one another, we have, for example,
\[
  \sum_a\hat l_{ay}\sum_a\hat l_{az}
  -\sum_a\hat l_{az}\sum_a\hat l_{ay}
  =\sum_a(\hat l_{ay}\hat l_{az}-\hat l_{az}\hat l_{ay})
  =i\sum_a\hat l_{ax}.
\]
Therefore
\begin{equation}
  \{\hat L_y,\hat L_z\}=i\hat L_x,\qquad
  \{\hat L_z,\hat L_x\}=i\hat L_y,\qquad
  \{\hat L_x,\hat L_y\}=i\hat L_z.
  \tag{26.8}
\end{equation}
Relations (26.8) show that the three components of angular momentum cannot
have definite values simultaneously (except when all three components vanish
at once; see below). In this respect angular momentum differs essentially
from momentum, whose three components can be measured simultaneously.

From $\hat L_x$, $\hat L_y$, and $\hat L_z$ we form the operator for the square
of the magnitude of the angular-momentum vector:
\begin{equation}
  \hat{\mathbf L}^{2}=\hat L_x^2+\hat L_y^2+\hat L_z^2.
  \tag{26.9}
\end{equation}
This operator commutes with each of $\hat L_x$, $\hat L_y$, and $\hat L_z$:
\begin{equation}
  \{\hat{\mathbf L}^{2},\hat L_x\}=0,\qquad
  \{\hat{\mathbf L}^{2},\hat L_y\}=0,\qquad
  \{\hat{\mathbf L}^{2},\hat L_z\}=0.
  \tag{26.10}
\end{equation}
Indeed, using (26.8), we find, for example,
\[
  \begin{aligned}
  \{\hat L_x^2,\hat L_z\}
    &=\hat L_x\{\hat L_x,\hat L_z\}
      +\{\hat L_x,\hat L_z\}\hat L_x
      =-i(\hat L_x\hat L_y+\hat L_y\hat L_x),\\
  \{\hat L_y^2,\hat L_z\}
    &=i(\hat L_x\hat L_y+\hat L_y\hat L_x),\\
  \{\hat L_z^2,\hat L_z\}&=0.
  \end{aligned}
\]
\SourcePage{113}{116}
Adding these equalities gives the last relation in (26.10).

Physically, (26.10) means that the square of the angular momentum (that is,
its magnitude) can have a definite value simultaneously with one of its
components.

Instead of $\hat L_x$ and $\hat L_y$, it is often convenient to use their
complex combinations
\begin{equation}
  \hat L_+=\hat L_x+i\hat L_y,\qquad
  \hat L_-=\hat L_x-i\hat L_y.
  \tag{26.11}
\end{equation}
Direct calculation from (26.8) readily gives the following commutation rules
for these combinations:
\begin{equation}
  \{\hat L_+,\hat L_-\}=2\hat L_z,\qquad
  \{\hat L_z,\hat L_+\}=\hat L_+,
  \qquad \{\hat L_z,\hat L_-\}=-\hat L_-.
  \tag{26.12}
\end{equation}
It is also easy to verify that
\begin{equation}
  \hat{\mathbf L}^{2}
  =\hat L_+\hat L_-+\hat L_z^2-\hat L_z
  =\hat L_-\hat L_++\hat L_z^2+\hat L_z.
  \tag{26.13}
\end{equation}

Finally, we give the frequently used expressions for the angular-momentum
operator of one particle in spherical coordinates. Introducing these
coordinates through the usual relations
\[
  x=r\sin\theta\cos\varphi,\qquad
  y=r\sin\theta\sin\varphi,\qquad
  z=r\cos\theta,
\]
a straightforward calculation gives
\begin{equation}
  \hat l_z=-i\frac{\partial}{\partial\varphi},
  \tag{26.14}
\end{equation}
\begin{equation}
  \hat l_\pm=e^{\pm i\varphi}
  \left(\pm\frac{\partial}{\partial\theta}
  +i\cot\theta\frac{\partial}{\partial\varphi}\right).
  \tag{26.15}
\end{equation}
Substitution in (26.13) gives the square of the particle's angular-momentum
operator in the form
\begin{equation}
  \hat{\mathbf l}^{2}=-\left[
  \frac{1}{\sin^2\theta}\frac{\partial^2}{\partial\varphi^2}
  +\frac{1}{\sin\theta}\frac{\partial}{\partial\theta}
  \left(\sin\theta\frac{\partial}{\partial\theta}\right)\right].
  \tag{26.16}
\end{equation}
Notice that, apart from a factor, this is the angular part of the Laplace
operator.
